\documentclass[11pt,a4paper,reqno]{amsart}

\usepackage[utf8]{inputenc}
\usepackage[T1]{fontenc}
\usepackage{amsmath,amssymb,amsfonts,amsthm,mathtools}
\usepackage{mathrsfs}
\usepackage{microtype}
\usepackage{booktabs}
\usepackage{tabularx}
\usepackage[dvipsnames]{xcolor}
\usepackage{hyperref}
\usepackage{geometry}
\usepackage{enumitem}
\usepackage{cite}

\geometry{
  top=3cm,
  bottom=3cm,
  left=2.8cm,
  right=2.8cm,
  headheight=14pt
}

\hypersetup{
  colorlinks=true,
  linkcolor=NavyBlue,
  citecolor=Mahogany,
  urlcolor=TealBlue,
  pdfauthor={Reinaldo M. Silva-Filho},
  pdftitle={Functional Realizations of Matrices and Hypertensors: Emergent Invariants and Geometric Measures}
}

% --- Theorem Environments ---
\newtheorem{theorem}{Theorem}[section]
\newtheorem{lemma}[theorem]{Lemma}
\newtheorem{proposition}[theorem]{Proposition}
\newtheorem{corollary}[theorem]{Corollary}

\theoremstyle{definition}
\newtheorem{definition}[theorem]{Definition}
\newtheorem{example}[theorem]{Example}
\newtheorem{remark}[theorem]{Remark}
\newtheorem{construction}[theorem]{Construction}

\numberwithin{equation}{section}

% --- Custom Notations ---
\providecommand{\R}{\mathbb{R}}
\newcommand{\C}{\mathbb{C}}
\newcommand{\N}{\mathbb{N}}
\newcommand{\Z}{\mathbb{Z}}
\providecommand{\Sph}{\mathbb{S}}
\newcommand{\Tor}{\mathbb{T}}
\newcommand{\Id}{\mathrm{Id}}
\newcommand{\Tr}{\operatorname{Tr}}
\newcommand{\rank}{\operatorname{rank}}
\newcommand{\supp}{\operatorname{supp}}
\newcommand{\Lip}{\operatorname{Lip}}
\providecommand{\BV}{\operatorname{BV}}
\providecommand{\TV}{\operatorname{TV}}
\newcommand{\Per}{\operatorname{Per}}
\providecommand{\cutnorm}[1]{\|#1\|_{\square}}
\newcommand{\norm}[1]{\left\|#1\right\|}
\newcommand{\inner}[2]{\left\langle #1, #2 \right\rangle}
\newcommand{\dif}{\,\mathrm{d}}
\newcommand{\abs}[1]{\left|#1\right|}

\title[Functional Realizations of Matrices and Hypertensors]{Functional Realizations of Matrices and Hypertensors:\\ Emergent Invariants and Geometric Measures}

\author{Reinaldo M. Silva-Filho}
\address{Universidade Federal de Lavras (UFLA), Departamento de Estat\'istica e Experimenta\c{c}\~ao Agropecu\'aria (PPGEE/DES)}
\email{reinaldo.filho1@estudante.ufla.br}
\date{\today}


\DeclareMathOperator*{\argmin}{argmin}
\providecommand{\II}{\mathrm{II}}

\begin{document}

\begin{abstract}
Finite-dimensional multilinear algebra treats matrices and higher-order tensors as elements of discrete linear spaces, extracting intrinsic invariants such as eigenvalues, singular values, operator rank, and trace polynomials. However, pure algebraic representations fundamentally obscure spatial locality, topological connectivity, and geometric regularity due to basis-permutation invariance in coordinates. In this treatise, we introduce a unified mathematical framework termed \emph{Functional Realization Operators} $\Phi \colon \mathcal{T}^k(V) \to \mathcal{F}(\Omega)$, mapping discrete algebraic arrays into function spaces over continuous domains $\Omega$. We construct a systematic taxonomy of continuous, piecewise continuous, step-like, and operator-valued realization families.

We prove, or recall with references, that embedding matrices and tensors into function spaces produces quantities that are not functions of the spectrum or of other permutation-invariant data: (i) Dirichlet energies and Sobolev $H^s(\Omega)$ norms, which differ between permutation-conjugate matrices; (ii) total variation ($\BV$) norms and the Hausdorff measures of level sets via the coarea formula; (iii) Morse indices of critical points and the Euler characteristic of spheres; and (iv) the cut-norm topology of graphons, which compares matrices of different sizes and has convergent subsequences as $n \to \infty$. These quantities are explicit functions of the entries; their interest lies in the spatial information they record.

We then discuss five application areas: (1) a Sobolev-regularized training penalty, proposed as a hypothesis for empirical study, and a graphon-based width-scaling protocol; (2) the Kac--Rice complexity of the isotropic spherical tensor landscape, which grows like $e^{\frac{n}{2}\log(d-1)}$ for degree $d \ge 3$ (Auffinger--Ben~Arous--\v{C}ern\'y), in contrast with the $2n$ critical points of the quadratic case; (3) size-independent comparison of large networks via the cut distance; (4) a well-posed variational formulation of the best rank-one tensor approximation (the rank-$r \ge 2$ ill-posedness of De Silva--Lim is not removed); and (5) the geometric (coarea) interpretation of total variation used in imaging.
\end{abstract}

\keywords{Matrix algebra, higher-order tensors, functional realizations, Sobolev spaces, total variation, Coarea formula, Morse theory, graphons, cut norm, spin glasses, deep learning.}
\subjclass[2020]{15A69, 46E35, 26B30, 58E05, 05C80, 47G10, 82B44, 68T07}

\maketitle

\tableofcontents

\newpage

% =========================================================================
\section{Introduction and Motivation}
\label{sec:intro}
% =========================================================================

Classical linear and multilinear algebra views a real matrix $A \in \R^{m \times n}$ or a tensor $\mathcal{T} \in \R^{d_1 \times \dots \times d_k}$ through the lens of coordinate arrays and linear maps between finite-dimensional vector spaces. Within this paradigm, the dominant structural invariants are purely algebraic: the trace $\Tr(A)$, the determinant $\det(A)$, the characteristic polynomial $P_A(\lambda)$, the spectrum $\sigma(A) = \{\lambda_i(A)\}$, the singular value spectrum $\Sigma(A) = \{\sigma_i(A)\}$, and the matrix rank $\rank(A)$.

While undeniably foundational, this purely algebraic framework suffers from a significant limitation when matrices represent spatial, physical, network, or image-based structures: \emph{algebraic invariants are oblivious to metric and topological spatial layouts}. Specifically:
\begin{enumerate}[label=(\alph*)]
    \item \textbf{Permutation Invariance vs. Metric Proximity:} If $P \in \R^{n \times n}$ is a permutation matrix, the conjugated matrix $B = P A P^T$ possesses an identical spectrum, trace, determinant, and Frobenius norm:
    \begin{equation}
        \sigma(P A P^T) = \sigma(A), \quad \Tr(P A P^T) = \Tr(A), \quad \|P A P^T\|_F = \|A\|_F.
    \end{equation}
    Yet, in physical space, transposing two adjacent rows versus swapping two distant rows induces fundamentally different topological discontinuities, gradient concentrations, and oscillatory phenomena.
    \item \textbf{The Curse of Dimensional Rigidity ($n \to \infty$):} Two matrices of differing sizes, say $A_{100 \times 100}$ and $B_{1000 \times 1000}$, live in non-isomorphic vector spaces. Algebraic linear algebra provides no intrinsic metric to measure how closely $A$ and $B$ approximate the same underlying continuum limit object without ad-hoc downsampling or artificial padding.
    \item \textbf{Higher-Order Tensor Intractability:} For tensors $\mathcal{T} \in \R^{d_1 \times d_2 \times d_3}$ of order $k \ge 3$, algebraic invariants become analytically intractable: tensor rank determination is NP-hard over $\R$ and $\C$, low-rank approximations can fail to exist due to non-closed orbits (border rank pathologies), and multilinear spectral theory yields complicated varieties rather than discrete eigenspaces.
\end{enumerate}

To transcend these limitations, one must bridge the gap between multilinear algebra and \emph{infinite-dimensional functional analysis, geometric measure theory, and differential topology}. 

\subsection*{Core Thesis and Contributions}
This paper demonstrates that by formally mapping matrices and tensors into function spaces via \emph{Functional Realization Operators} $\Phi \colon \mathcal{T}^k(V) \to \mathcal{F}(\Omega)$, we unlock a spectrum of geometric and analytical invariants that:
\begin{enumerate}[label=\arabic*.]
    \item Discriminate between spectrally indistinguishable arrays according to spatial energy, harmonic smoothness, and Sobolev scale $H^s(\Omega)$;
    \item Measure internal perimeter and level-set geometry through functions of bounded variation ($\BV$) and the geometric Coarea formula;
    \item Unveil the non-convex optimization topology through Morse critical point indices on compact manifolds ($\Sph^{n-1}$);
    \item Establish continuous limits of dense discrete arrays via the cut-norm topology and graphons ($L^\infty \to L^1$ operator duality);
    \item Connect these invariants to five application areas (machine learning, spin glass physics, large networks, tensor approximation, and imaging), distinguishing proved statements from proposals and conjectures.
\end{enumerate}

% =========================================================================
\section{Axiomatic Framework: Functional Realization Operators}
\label{sec:framework}
% =========================================================================

Let $V_1, \dots, V_k$ be real vector spaces of finite dimensions $d_1, \dots, d_k$ equipped with standard inner products $\inner{\cdot}{\cdot}_{V_i}$, and let $\mathcal{T}^k(V_1, \dots, V_k) \cong V_1 \otimes \dots \otimes V_k$ denote the $k$-th order tensor product space. When $V_i = \R^{d_i}$ with canonical bases $\{e_{j}^{(i)}\}$, any tensor $\mathcal{T} \in \mathcal{T}^k(V_1, \dots, V_k)$ is uniquely identified with its coordinate multi-array $(T_{j_1 \dots j_k})_{1 \le j_r \le d_r}$. For $k=2$, this reduces to the classical matrix space $\R^{d_1 \times d_2}$.

Let $\Omega \subseteq \R^d$ be a measurable domain endowed with the Lebesgue measure $\dif x$, or a smooth compact Riemannian manifold $(M, g)$ with volume form $\dif V_g$. We denote by $\mathcal{F}(\Omega)$ a topological vector space of functions over $\Omega$ (e.g., $C^0(\Omega)$, $L^p(\Omega)$, $W^{s,p}(\Omega)$, or the space of Radon measures $\mathcal{M}(\Omega)$).

\begin{definition}[Functional Realization Operator]
\label{def:realization_op}
A \emph{Functional Realization Operator} is a mapping
\begin{equation}
    \Phi \colon \mathcal{T}^k(V_1, \dots, V_k) \longrightarrow \mathcal{F}(\Omega)
\end{equation}
which associates to each discrete tensor $\mathcal{T}$ a well-defined function $f_{\mathcal{T}} \coloneqq \Phi(\mathcal{T}) \in \mathcal{F}(\Omega)$.
\end{definition}

\begin{definition}[Structural Properties of Realizations]
Let $\Phi \colon \mathcal{T}^k \to \mathcal{F}(\Omega)$ be a functional realization. We categorize $\Phi$ according to the following properties:
\begin{enumerate}[label=(\roman*)]
    \item \textbf{Linearity:} $\Phi$ is linear if $\Phi(\alpha \mathcal{S} + \beta \mathcal{T}) = \alpha \Phi(\mathcal{S}) + \beta \Phi(\mathcal{T})$ for all $\alpha, \beta \in \R$ and $\mathcal{S}, \mathcal{T} \in \mathcal{T}^k$.
    \item \textbf{Injectivity (Faithfulness):} $\Phi$ is faithful if $\ker(\Phi) = \{0\}$, meaning no non-zero tensor vanishes identically as a continuous function.
    \item \textbf{Continuity (Boundedness):} If $\mathcal{T}^k$ is equipped with the Frobenius norm $\|\mathcal{T}\|_F = \sqrt{\sum T_{j_1 \dots j_k}^2}$ and $\mathcal{F}(\Omega)$ with norm $\|\cdot\|_{\mathcal{F}}$, $\Phi$ is bounded if there exists $C > 0$ such that
    \begin{equation}
        \|\Phi(\mathcal{T})\|_{\mathcal{F}} \le C \|\mathcal{T}\|_F, \quad \forall \mathcal{T} \in \mathcal{T}^k.
    \end{equation}
    \item \textbf{Symmetry Equivariance:} Let $G$ be a subgroup of the product of orthogonal groups $O(d_1) \times \dots \times O(d_k)$ (or permutation groups $\mathcal{S}_{d_1} \times \dots \times \mathcal{S}_{d_k}$) acting on $\mathcal{T}^k$, and let $G$ act on $\Omega$ via measurable transformations $\tau_g \colon \Omega \to \Omega$. $\Phi$ is $G$-equivariant if
    \begin{equation}
        \Phi(g \cdot \mathcal{T})(x) = \Phi(\mathcal{T})(\tau_{g^{-1}}(x)), \quad \forall x \in \Omega, \forall g \in G.
    \end{equation}
\end{enumerate}
\end{definition}

% -------------------------------------------------------------------------
\section{Comprehensive Taxonomy of Functional Realizations}
\label{sec:taxonomy}
% -------------------------------------------------------------------------

\subsection{Family I: Multilinear, Polynomial, and Spherical Realizations}
\label{subsec:fam1}

In this family, the matrix or tensor coefficients define the homogeneous algebraic form of a polynomial on Euclidean space or a compact sphere.

\begin{construction}[Quadratic Forms on $\R^n$ and $\Sph^{n-1}$]
Let $A \in \R^{n \times n}$. The quadratic realization $\Phi_{\mathrm{quad}}(A) \colon \R^n \to \R$ is defined by:
\begin{equation}
    \label{eq:quad_form}
    f_A(x) \coloneqq x^T A x = \sum_{i=1}^n \sum_{j=1}^n a_{ij} x_i x_j.
\end{equation}
Restricting $x$ to the unit sphere $\Sph^{n-1} = \{x \in \R^n \colon \|x\|_2 = 1\}$ yields a smooth function $f_A \in C^\infty(\Sph^{n-1})$.
\end{construction}

\begin{construction}[Higher-Order Homogeneous Tensor Forms]
For a $d$-th order symmetric tensor $\mathcal{T} \in \operatorname{Sym}^d(\R^n)$, the homogeneous realization $\Phi_{\mathrm{hom}}(\mathcal{T}) \colon \Sph^{n-1} \to \R$ is:
\begin{equation}
    \label{eq:tensor_form}
    f_{\mathcal{T}}(x) \coloneqq \mathcal{T}(x, x, \dots, x) = \sum_{i_1, \dots, i_d = 1}^n \mathcal{T}_{i_1 \dots i_d} x_{i_1} \dots x_{i_d}.
\end{equation}
\end{construction}

\begin{proposition}[Spectral and Critical Correspondence]
Let $A \in \operatorname{Sym}_n(\R)$ be symmetric with eigendecomposition $A = V \Lambda V^T$, where $\Lambda = \operatorname{diag}(\lambda_1, \dots, \lambda_n)$ with $\lambda_1 \le \lambda_2 \le \dots \le \lambda_n$. Then:
\begin{enumerate}[label=(\roman*)]
    \item The global extrema of $f_A|_{\Sph^{n-1}}$ satisfy:
    \begin{equation}
        \min_{x \in \Sph^{n-1}} f_A(x) = \lambda_1, \quad \max_{x \in \Sph^{n-1}} f_A(x) = \lambda_n.
    \end{equation}
    \item The critical points of the function $f_A \colon \Sph^{n-1} \to \R$ are precisely the normalized eigenvectors $v_i$ of $A$, and the critical values are the corresponding eigenvalues $\lambda_i$.
\end{enumerate}
\end{proposition}
\begin{proof}
Consider the Lagrangian $\mathcal{L}(x, \mu) = x^T A x - \mu (x^T x - 1)$ on $\R^n \times \R$. The Euler-Lagrange criticality condition $\nabla_x \mathcal{L} = 2Ax - 2\mu x = 0$ immediately yields the eigenvalue equation $Ax = \mu x$. Constraining $\|x\|_2 = 1$ gives $f_A(x) = x^T A x = \mu x^T x = \mu$. Thus, the critical values are precisely the spectrum $\sigma(A)$, and the extrema coincide with the Rayleigh quotient bounds.
\end{proof}

\subsection{Family II: Step-like, Partition-Based, and Graphon Realizations}
\label{subsec:fam2}

Here, the discrete array parametrizes a piecewise-constant function over a partitioned measure space. This construction forms the foundation of modern dense graph limit theory (Graphons) and digital imaging.

\begin{construction}[Graphon / Unit Square Step Realization]
Let $A \in \R^{n \times n}$ and consider the unit square $\Omega = [0, 1]^2$ with Lebesgue measure. Partition $[0, 1]$ into $n$ equal left-closed, right-open intervals:
\begin{equation}
    I_i \coloneqq \left[\frac{i-1}{n}, \frac{i}{n}\right), \quad i=1, \dots, n-1, \quad \text{and} \quad I_n \coloneqq \left[\frac{n-1}{n}, 1\right].
\end{equation}
The graphon step realization $\Phi_{\mathrm{step}}(A) \colon [0, 1]^2 \to \R$ is defined by:
\begin{equation}
    \label{eq:graphon_step}
    W_A(x, y) \coloneqq \sum_{i=1}^n \sum_{j=1}^n a_{ij} \mathbf{1}_{I_i}(x) \mathbf{1}_{I_j}(y),
\end{equation}
where $\mathbf{1}_E$ denotes the characteristic function of set $E$.
\end{construction}

\begin{construction}[Multidimensional Pixel/Voxel Fields]
Let $A \in \R^{n_1 \times \dots \times n_d}$ be a $d$-dimensional tensor. Partition $\Omega = \prod_{\alpha=1}^d [0, L_\alpha]$ into hypercubes $Q_{\mathbf{i}} = \prod_{\alpha=1}^d [(\mathbf{i}_\alpha-1) h_\alpha, \mathbf{i}_\alpha h_\alpha)$. The realization is:
\begin{equation}
    f_A(x) = \sum_{\mathbf{i}} A_{\mathbf{i}} \mathbf{1}_{Q_{\mathbf{i}}}(x).
\end{equation}
Notice that $f_A \notin C^0(\Omega)$; rather, $f_A \in L^\infty(\Omega) \cap \BV(\Omega)$.
\end{construction}

\begin{construction}[Polyhedral Indicator Functions and Hyperplane Arrangements]
Let $A \in \R^{m \times n}$ and $b \in \R^m$. The matrix $A$ defines a system of $m$ affine half-spaces $H_i^- = \{x \in \R^n \colon \inner{a_i}{x} \le b_i\}$, where $a_i$ is the $i$-th row of $A$. The polyhedral indicator realization is:
\begin{equation}
    \Phi_{\mathrm{poly}}(A, b)(x) \coloneqq \mathbf{1}_{\{x \in \R^n \colon Ax \le b\}}(x) = \prod_{i=1}^m \Theta(b_i - \inner{a_i}{x}),
\end{equation}
where $\Theta(t) = \mathbf{1}_{[0, \infty)}(t)$ is the Heaviside step function.
\end{construction}

\subsection{Family III: Interpolation, Basis Expansions, and Sobolev Fields}
\label{subsec:fam3}

In this family, matrix elements act as generalized Fourier coefficients, spline control points, or localized radial basis weights.

\begin{construction}[Harmonic / Dual Fourier Realization]
Let $A \in \C^{n \times n}$. Let $\Tor^2 = [0, 2\pi)^2$ be the 2-torus equipped with the normalized Haar probability measure $\dif\mu(x) = \frac{\dif x_1 \dif x_2}{(2\pi)^2}$. We associate $A$ with the trigonometric polynomial:
\begin{equation}
    \label{eq:fourier_realization}
    f_A(x_1, x_2) \coloneqq \sum_{k_1, k_2 = 1}^n a_{k_1 k_2} e^{i (k_1 x_1 + k_2 x_2)}.
\end{equation}
By Parseval's identity on $\{e^{i(k_1 x_1 + k_2 x_2)}\}$, the $L^2(\Tor^2, \dif\mu)$ norm satisfies $\|f_A\|_{L^2(\Tor^2)} = \|A\|_F$.
\end{construction}

\begin{construction}[Tensor Product Splines and B\'ezier Surfaces]
Let $B_{i, p}(u)$ denote the $i$-th B-spline basis function of degree $p$ on a knot vector $\Xi$, and $B_{j,q}(v)$ those of degree $q$ on a knot vector $H$. For an array of control points $P \in \R^{(m+1) \times (n+1) \times 3}$, indexed by $0 \le i \le m$ and $0 \le j \le n$:
\begin{equation}
    S(u, v) = \sum_{i=0}^m \sum_{j=0}^n B_{i, p}(u) B_{j, q}(v) P_{ij}.
\end{equation}
The surface is $C^{p-1}$ in $u$ and $C^{q-1}$ in $v$ at simple knots (a knot of multiplicity $\mu$ lowers the regularity there to $C^{p-\mu}$). The control points determine the shape through the convex hull property and the curvature of $S$.
\end{construction}

\subsection{Family IV: Integral Operators and Continuous Matrix Limits}
\label{subsec:fam4}

\begin{construction}[Fredholm Integral Kernel Operators]
Every matrix $A \in \R^{n \times n}$ can be identified with the integral operator $T_A \colon L^2([0, 1]) \to L^2([0, 1])$:
\begin{equation}
    (T_A u)(x) \coloneqq \int_0^1 W_A(x, y) u(y) \dif y,
\end{equation}
where $W_A$ is the step realization given in \eqref{eq:graphon_step}. $T_A$ is a Hilbert-Schmidt operator with Hilbert-Schmidt norm $\|T_A\|_{\mathrm{HS}} = \frac{1}{n} \|A\|_F$.
\end{construction}

% =========================================================================
\section{Emerging Invariants: What Matrix Algebra Cannot See}
\label{sec:emerging_invariants}
% =========================================================================

\subsection{Spatiality vs. Spectrum: Dirichlet Energies and Sobolev Regularity}
\label{subsec:sobolev}

Let $A \in \R^{n \times n}$. In linear algebra, the permutation matrix $P_\pi$ corresponding to a permutation $\pi \in \mathcal{S}_n$ preserves all spectral invariants:
\begin{equation}
    \sigma(P_\pi A P_\pi^T) = \sigma(A), \quad \Tr(P_\pi A P_\pi^T) = \Tr(A), \quad \|P_\pi A P_\pi^T\|_F = \|A\|_F.
\end{equation}

Now consider the harmonic realization on $(\Tor^2, \dif\mu)$ given by \eqref{eq:fourier_realization}:
\begin{equation}
    f_A(x_1, x_2) = \sum_{k_1, k_2 = 1}^n a_{k_1 k_2} e^{i (k_1 x_1 + k_2 x_2)}.
\end{equation}

The homogeneous Dirichlet energy under normalized measure $\dif\mu(x) = \frac{\dif x_1 \dif x_2}{(2\pi)^2}$ is:
\begin{equation}
    \mathcal{E}(f) \coloneqq \|\nabla f\|_{L^2(\Tor^2, \dif\mu)}^2 = \int_{\Tor^2} \|\nabla f(x_1, x_2)\|_2^2 \dif\mu(x).
\end{equation}

\begin{theorem}[Spectral Blindness to Spatial Energy]
\label{thm:dirichlet_blindness}
For any matrix $A \in \R^{n \times n}$, the Dirichlet energy of the realization $f_A$ is given explicitly by:
\begin{equation}
    \label{eq:dirichlet_energy_formula}
    \mathcal{E}(f_A) = \sum_{k_1=1}^n \sum_{k_2=1}^n (k_1^2 + k_2^2) |a_{k_1 k_2}|^2.
\end{equation}
Moreover, there exist non-zero symmetric matrices $A \in \operatorname{Sym}_n(\R)$ and orthogonal permutation matrices $P_1, P_2 \in \mathcal{S}_n$ such that $A_1 = P_1 A P_1^T$ and $A_2 = P_2 A P_2^T$ satisfy:
\begin{equation}
    \sigma(A_1) = \sigma(A_2) \quad \text{and} \quad \|A_1\|_F = \|A_2\|_F = \|A\|_F,
\end{equation}
yet the ratio of their functional Dirichlet energies scales as:
\begin{equation}
    \frac{\mathcal{E}(f_{A_2})}{\mathcal{E}(f_{A_1})} = \Theta(n^2).
\end{equation}
\end{theorem}

\begin{remark}[Isotropic Invariance vs Anisotropic Discrepancy]
For isotropic matrices such as $A = c I_n$ (scalar multiples of the identity) or $A = c J_n$ (the all-ones matrix), any permutation matrix satisfies $P A P^T = A$, yielding an energy ratio $\mathcal{E}(f_{A_2})/\mathcal{E}(f_{A_1}) = 1$. The maximal quadratic amplification $\Theta(n^2)$ is achieved by anisotropic and localized configurations, such as projectors and rank-one forms.
\end{remark}

\begin{proof}
Direct differentiation of $f_A$ yields:
\begin{equation}
    \partial_{x_1} f_A = \sum_{k_1, k_2=1}^n (i k_1) a_{k_1 k_2} e^{i(k_1 x_1 + k_2 x_2)}, \quad
    \partial_{x_2} f_A = \sum_{k_1, k_2=1}^n (i k_2) a_{k_1 k_2} e^{i(k_1 x_1 + k_2 x_2)}.
\end{equation}
Since this is a finite trigonometric polynomial, termwise differentiation and the interchange of limits are unconditionally justified. By Parseval's identity on the orthogonal character system $\{e^{i(k_1 x_1 + k_2 x_2)}\}$:
\begin{equation}
    \|\nabla f_A\|_{L^2}^2 = \|\partial_{x_1} f_A\|_{L^2}^2 + \|\partial_{x_2} f_A\|_{L^2}^2 = \sum_{k_1, k_2=1}^n (k_1^2 + k_2^2) |a_{k_1 k_2}|^2.
\end{equation}
To prove the ratio bound, consider a rank-1 matrix $A = u u^T$ where $u \in \R^n$ has a single non-zero component $u_1 = 1$. Let $\pi_1 = \mathrm{id}$ so that $A_1 = E_{1,1}$, which places the entire spectral mass at the base frequency index $(1, 1)$. Then:
\begin{equation}
    \mathcal{E}(f_{A_1}) = (1^2 + 1^2) \cdot 1^2 = 2.
\end{equation}
Now let $\pi_2$ swap index $1$ and $n$, so $A_2 = P_{\pi_2} A P_{\pi_2}^T = E_{n, n}$, concentrating the spectral mass at maximal frequency $(n, n)$. Then:
\begin{equation}
    \mathcal{E}(f_{A_2}) = (n^2 + n^2) \cdot 1^2 = 2n^2.
\end{equation}
Both $A_1$ and $A_2$ possess an identical spectrum $\sigma = \{1, 0, \dots, 0\}$, identical rank 1, and identical Frobenius norm $\|A_1\|_F = \|A_2\|_F = 1$. However:
\begin{equation}
    \frac{\mathcal{E}(f_{A_2})}{\mathcal{E}(f_{A_1})} = \frac{2n^2}{2} = n^2 = \Theta(n^2).
\end{equation}
This completes the proof.
\end{proof}

\begin{example}[The Checkerboard Matrix Pathology]
Consider the checkerboard array $A \in \R^{n \times n}$ given by $a_{ij} = (-1)^{i+j}$. In discrete linear algebra, $A = v v^T$ with $v = (1, -1, 1, -1, \dots)^T$ is an elementary rank-1 matrix with spectrum $\sigma(A) = \{n, 0, \dots, 0\}$. Pure algebra classifies $A$ as structurally simple. However, under the harmonic realization, every entry has $|a_{k_1 k_2}|^2 = 1$, yielding by \eqref{eq:dirichlet_energy_formula} a rapidly growing Dirichlet energy:
\begin{equation}
    \mathcal{E}(f_A) = \sum_{k_1=1}^n \sum_{k_2=1}^n (k_1^2 + k_2^2) \cdot 1 = 2n \sum_{k=1}^n k^2 = 2n \cdot \frac{n(n+1)(2n+1)}{6} = \frac{2}{3} n^4 + \mathcal{O}(n^3) = \Theta(n^4).
\end{equation}
Since $|a_{k_1 k_2}| \le 1$ for every entrywise-bounded matrix, \eqref{eq:dirichlet_energy_formula} shows that $A$ attains the largest possible Dirichlet energy, $\Theta(n^4)$, among all matrices with entries in $[-1,1]$. The same value is attained by \emph{every} $\pm 1$ matrix, including the all-ones matrix $J_n$. The harmonic realization treats the indices $(k_1,k_2)$ as frequencies, so its energy measures how mass is distributed over frequency indices, not the spatial sign oscillation of the checkerboard. Spatial roughness in the latter sense is captured by the step realization and its total variation (Theorem~\ref{thm:tv_matrix}), which is $\Theta(n)$ for the checkerboard and $0$ for $J_n$.
\end{example}

\subsection{Geometric Measure Theory: Bounded Variation and the Coarea Formula}
\label{subsec:coarea}

We now investigate the step realization $\Phi_{\mathrm{step}}(A)$ on $\Omega = (0, 1)^2$. Because $W_A$ is discontinuous across grid edges, its classical gradient does not exist. However, $W_A$ belongs to the space of \emph{Functions of Bounded Variation} $\BV(\Omega)$.

\begin{definition}[Total Variation in $\BV(\Omega)$]
Let $\Omega \subset \R^2$ be an open domain. For $f \in L^1(\Omega)$, the \emph{Total Variation} of $f$ is defined duality-wise by:
\begin{equation}
    \TV(f) \coloneqq \sup \left\{ \int_\Omega f(x) \operatorname{div} \varphi(x) \dif x \colon \varphi \in C_c^1(\Omega; \R^2), \ \|\varphi\|_{L^\infty} \le 1 \right\}.
\end{equation}
\end{definition}

\begin{theorem}[Total Variation of the Matrix Step Realization]
\label{thm:tv_matrix}
Let $A \in \R^{n \times n}$ and let $W_A \in \BV((0,1)^2)$ be its step realization defined on the open unit square $(0, 1)^2$. Then the distributional gradient $D W_A$ is a 1-dimensional Radon measure concentrated on the grid lines $\mathcal{G}_n = \bigcup_{i=1}^{n-1} (\{i/n\} \times [0,1]) \cup \bigcup_{j=1}^{n-1} ([0,1] \times \{j/n\})$, and its total variation is given by:
\begin{equation}
    \label{eq:tv_formula}
    \TV(W_A) = \frac{1}{n} \sum_{i=1}^{n-1} \sum_{j=1}^n |a_{i+1, j} - a_{i, j}| + \frac{1}{n} \sum_{i=1}^n \sum_{j=1}^{n-1} |a_{i, j+1} - a_{i, j}|.
\end{equation}
\end{theorem}

\begin{proof}
Let $\varphi = (\varphi_1, \varphi_2) \in C_c^1((0,1)^2; \R^2)$ with $\|\varphi\|_{L^\infty} \le 1$. By Green-Gauss integration by parts over each sub-square $Q_{ij} = (\frac{i-1}{n}, \frac{i}{n}) \times (\frac{j-1}{n}, \frac{j}{n})$:
\begin{equation}
    \int_{(0,1)^2} W_A \operatorname{div} \varphi \dif x \dif y = \sum_{i,j=1}^n a_{ij} \int_{Q_{ij}} \left( \frac{\partial \varphi_1}{\partial x} + \frac{\partial \varphi_2}{\partial y} \right) \dif x \dif y.
\end{equation}
Applying the divergence theorem on each cell gives:
\begin{equation}
    \int_{Q_{ij}} \operatorname{div}\varphi \dif x \dif y = \int_{\partial Q_{ij}} \varphi \cdot \nu \dif \mathcal{H}^1,
\end{equation}
where $\nu$ is the outward unit normal and $\mathcal{H}^1$ is the 1-dimensional Hausdorff measure. 

At the internal vertical boundary separating $Q_{ij}$ and $Q_{i+1, j}$, denoted $\Gamma_{i, j}^v = \{i/n\} \times (\frac{j-1}{n}, \frac{j}{n})$, the outward normal from $Q_{ij}$ is $(1, 0)$ and from $Q_{i+1, j}$ is $(-1, 0)$. Likewise, the horizontal interface between $Q_{ij}$ and $Q_{i,j+1}$ is $\Gamma_{i, j}^h = (\frac{i-1}{n}, \frac{i}{n}) \times \{j/n\}$, with outward normals $(0,1)$ from $Q_{ij}$ and $(0,-1)$ from $Q_{i,j+1}$. The test field vanishes near $\partial(0,1)^2$, so the outer boundary contributes nothing. Summing across all interfaces:
\begin{align}
    \int_{(0,1)^2} W_A \operatorname{div} \varphi \dif x \dif y 
    &= \sum_{i=1}^{n-1} \sum_{j=1}^n (a_{ij} - a_{i+1, j}) \int_{\Gamma_{i,j}^v} \varphi_1\left(\frac{i}{n}, y\right) \dif y \notag \\
    &\quad + \sum_{i=1}^n \sum_{j=1}^{n-1} (a_{ij} - a_{i, j+1}) \int_{\Gamma_{i,j}^h} \varphi_2\left(x, \frac{j}{n}\right) \dif x.
\end{align}
Taking the supremum over test vector fields with support concentrated in disjoint tubular neighborhoods of the interface segments allows us to saturate the sign of each difference independently on each segment of length $1/n$:
\begin{equation}
    \sup_{\|\varphi\|\le 1} \int_{(0,1)^2} W_A \operatorname{div} \varphi \dif x \dif y = \sum_{i=1}^{n-1} \sum_{j=1}^n |a_{i+1, j} - a_{i,j}| \cdot \frac{1}{n} + \sum_{i=1}^n \sum_{j=1}^{n-1} |a_{i, j+1} - a_{i, j}| \cdot \frac{1}{n}.
\end{equation}
This establishes \eqref{eq:tv_formula}.
\end{proof}

\begin{theorem}[Geometric Coarea Linkage: Hausdorff Measure of Matrix Level Curves]
\label{thm:coarea_linkage}
For any $t \in \R$, define the super-level set $E_t \coloneqq \{(x, y) \in (0, 1)^2 \colon W_A(x, y) > t\}$. Then $E_t$ is a set of finite perimeter in $(0, 1)^2$, and its perimeter $\Per(E_t; (0,1)^2) = \mathcal{H}^1(\partial^* E_t \cap (0,1)^2)$ satisfies the \emph{coarea formula}:
\begin{equation}
    \label{eq:coarea}
    \TV(W_A) = \int_{-\infty}^\infty \mathcal{H}^1(\partial^* E_t \cap (0,1)^2) \dif t,
\end{equation}
where $\partial^* E_t$ is the reduced boundary of $E_t$ in the sense of De Giorgi.
\end{theorem}
\begin{proof}
Since $W_A \in \BV((0,1)^2)$ by Theorem~\ref{thm:tv_matrix}, the coarea formula for functions of bounded variation of Fleming and Rishel~\cite{fleming1960integral} gives $\TV(W_A) = \int_{\R} \Per(E_t;(0,1)^2) \dif t$, with $E_t$ of finite perimeter for almost every $t$; see also Federer~\cite{federer1969geometric}, Ambrosio, Fusco, and Pallara~\cite{ambrosio2000functions}, and Evans and Gariepy~\cite{evans2015measure}. De Giorgi's structure theorem identifies $\Per(E_t;(0,1)^2)$ with $\mathcal{H}^1(\partial^* E_t \cap (0,1)^2)$. Here $W_A$ takes finitely many values, so $E_t$ is a finite union of grid cells and has finite perimeter for every $t$.
\end{proof}

\subsection{Differential Topology: Morse Theory and Landscape Homology}
\label{subsec:morse}

We now examine the polynomial realization on the sphere $\Sph^{n-1}$, defined in Section~\ref{subsec:fam1}:
\begin{equation}
    f_A(x) = x^T A x, \quad x \in \Sph^{n-1}.
\end{equation}
Assume that $A \in \operatorname{Sym}_n(\R)$ is symmetric.

\begin{theorem}[Morse Spectrum of Quadratic Realizations]
\label{thm:morse_matrix}
Let $A \in \operatorname{Sym}_n(\R)$ have strictly distinct eigenvalues:
\begin{equation}
    \lambda_1 < \lambda_2 < \dots < \lambda_n.
\end{equation}
Then $f_A \colon \Sph^{n-1} \to \R$ is a strictly non-degenerate Morse function. Moreover:
\begin{enumerate}[label=(\roman*)]
    \item $f_A$ has exactly $2n$ isolated critical points, occurring in antipodal pairs:
    \begin{equation}
        \operatorname{Crit}(f_A) = \{\pm v_1, \pm v_2, \dots, \pm v_n\},
    \end{equation}
    where $v_i$ is the normalized eigenvector corresponding to $\lambda_i$.
    \item The Morse index of the critical points $\pm v_i$ is given precisely by:
    \begin{equation}
        \gamma(\pm v_i) = i - 1.
    \end{equation}
    \item The Morse polynomial $P_t(f_A) \coloneqq \sum_{p \in \operatorname{Crit}(f_A)} t^{\gamma(p)}$ satisfies:
    \begin{equation}
        P_t(f_A) = 2 \sum_{i=1}^n t^{i-1} = 2 (1 + t + t^2 + \dots + t^{n-1}).
    \end{equation}
    Evaluating at $t = -1$ rigorously recovers the topological Euler--Poincar\'e characteristic:
    \begin{equation}
        \chi(\Sph^{n-1}) = P_{-1}(f_A) = 1 - (-1)^n = \begin{cases} 0, & n \text{ even} \ (\Sph^{n-1} \text{ odd}), \\ 2, & n \text{ odd} \ (\Sph^{n-1} \text{ even}). \end{cases}
    \end{equation}
\end{enumerate}
\end{theorem}

\begin{proof}
Let $x \in \Sph^{n-1}$. The Riemannian gradient on the sphere with respect to the standard round metric is the projection of the ambient gradient onto $T_x \Sph^{n-1} = \{u \in \R^n \colon \inner{u}{x} = 0\}$:
\begin{equation}
    \nabla_{\Sph^{n-1}} f_A(x) = 2 A x - 2(x^T A x) x.
\end{equation}
Thus, $\nabla_{\Sph^{n-1}} f_A(x) = 0$ if and only if $Ax = (x^T Ax) x$. Since $\|x\|_2=1$, $x$ is an eigenvector with eigenvalue $\lambda = x^T Ax$. Because the eigenvalues are strictly distinct, the only unit critical points are $x = \pm v_i$ for $i=1, \dots, n$.

Next, we evaluate the Riemannian Hessian at $x = v_i$. Along a geodesic $\gamma(t)$ on $\Sph^{n-1}$ with $\gamma(0) = v_i$ and $\dot{\gamma}(0) = u \in T_{v_i} \Sph^{n-1}$, the geodesic curvature satisfies $\ddot{\gamma}(0) = -\|u\|_2^2 v_i$. Differentiating:
\begin{equation}
    \operatorname{Hess}_{\Sph^{n-1}} f_A(v_i)(u, u) = \left.\frac{\dif^2}{\dif t^2} f_A(\gamma(t))\right|_{t=0} = 2 u^T A u - 2 \lambda_i \|u\|_2^2.
\end{equation}
Expanding $u$ in the orthonormal eigenbasis $\{v_j\}_{j \neq i}$: $u = \sum_{j \ne i} c_j v_j$. Then:
\begin{equation}
    \operatorname{Hess}_{\Sph^{n-1}} f_A(v_i)(u, u) = 2 \sum_{j \ne i} (\lambda_j - \lambda_i) c_j^2.
\end{equation}
Since $\lambda_1 < \dots < \lambda_n$:
\begin{itemize}
    \item $(\lambda_j - \lambda_i) < 0$ for $j \in \{1, \dots, i-1\}$, producing exactly $i-1$ negative eigenvalues;
    \item $(\lambda_j - \lambda_i) > 0$ for $j \in \{i+1, \dots, n\}$, producing $(n-1) - (i-1) = n - i$ positive eigenvalues;
    \item there are no zero eigenvalues since $\lambda_j \ne \lambda_i$ for all $j \ne i$.
\end{itemize}
Thus, $\ker(\operatorname{Hess}_{\Sph^{n-1}} f_A(v_i)) = \{0\}$ and the Morse index is $\gamma(\pm v_i) = i - 1$. 

The Morse polynomial evaluates to $P_{-1}(f_A) = 2 \sum_{i=1}^n (-1)^{i-1} = 2 \cdot \frac{1 - (-1)^n}{2} = 1 - (-1)^n$. By classical Morse theory (see Milnor~\cite{milnor1963morse}), this alternating sum of Morse indices exactly recovers the Euler characteristic $\chi(\Sph^{n-1})$ obtained from singular homology $H_*(\Sph^{n-1})$.
\end{proof}

\begin{remark}[The Degenerate Morse-Bott Regime]
If $A$ has repeated eigenvalues, say $\lambda_k = \dots = \lambda_{k+m}$, the critical set along that eigenspace forms a continuous critical sub-manifold $\Sph^m$ rather than isolated points. In this setting, $f_A$ is a Morse-Bott function. The hypothesis of distinct eigenvalues in Theorem~\ref{thm:morse_matrix} is therefore essential and cannot be relaxed without changing the conclusion.
\end{remark}

\subsection{Asymptotic Graphons and the Cut-Norm Topology (\texorpdfstring{$n \to \infty$}{n -> infty})}
\label{subsec:graphon_cutnorm}

In the functional domain, both $W_A$ and $W_B$ reside in the universal Lebesgue space $\mathcal{W} \coloneqq \{W \in L^\infty([0, 1]^2) \colon W(x, y) = W(y, x)\}$.

\begin{definition}[The Cut Norm (Frieze--Kannan / Lov\'asz--Szegedy)]
For any kernel $W \in L^1([0, 1]^2)$, the \emph{Cut Norm} is defined by:
\begin{equation}
    \label{eq:cut_norm_def}
    \cutnorm{W} \coloneqq \sup_{S, T \subseteq [0, 1]} \left| \int_{S \times T} W(x, y) \dif x \dif y \right|,
\end{equation}
where the supremum is taken over all measurable subsets $S, T \subseteq [0, 1]$.
\end{definition}

\begin{theorem}[Equivalence with the $L^\infty \to L^1$ Operator Norm]
\label{thm:grothendieck_cutnorm}
The cut norm $\cutnorm{W}$ is equivalent to the operator norm of the integral operator $T_W \colon L^\infty([0,1]) \to L^1([0,1])$:
\begin{equation}
    \cutnorm{W} \le \|T_W\|_{L^\infty \to L^1} \le 4 \cutnorm{W},
\end{equation}
where $\|T_W\|_{L^\infty \to L^1} \coloneqq \sup_{\|u\|_\infty \le 1, \|v\|_\infty \le 1} \left| \iint_{[0,1]^2} W(x, y) u(x) v(y) \dif x \dif y \right|$.
\end{theorem}

\begin{proof}
For any indicator functions $u = \mathbf{1}_S$ and $v = \mathbf{1}_T$, we have $\inner{u}{T_W v} = \int_{S \times T} W(x,y) \dif x \dif y$. Decomposing any $u, v$ with $\|u\|_\infty \le 1, \|v\|_\infty \le 1$ into positive and negative parts $u = u_+ - u_-$ and $v = v_+ - v_-$:
\begin{align}
    \left|\int_{[0,1]^2} W u v \dif x \dif y\right| 
    &\le \left|\int W u_+ v_+\right| + \left|\int W u_+ v_-\right| + \left|\int W u_- v_+\right| + \left|\int W u_- v_-\right| \notag \\
    &\le 4 \cutnorm{W}.
\end{align}
Taking the supremum establishes the upper bound. The lower bound $\cutnorm{W} \le \|T_W\|_{L^\infty \to L^1}$ follows immediately by restricting $u = \mathbf{1}_S$ and $v = \mathbf{1}_T$.
\end{proof}

\begin{theorem}[Compactness of the Continuum Moduli Space]
\label{thm:graphon_compactness}
Let $\bar{\mathcal{W}}$ be the space of symmetric graphons $W \colon [0, 1]^2 \to [0, 1]$, let $\bar{\Pi}$ be the group of measure-preserving bijections of $[0,1]$, and let $\delta_\square(U, W) \coloneqq \inf_{\psi \in \bar\Pi} \cutnorm{U - W^\psi}$ with $W^\psi(x,y) = W(\psi(x),\psi(y))$. Identify $U$ and $W$ when $\delta_\square(U, W) = 0$ (weak isomorphism; this is coarser than lying in one $\bar\Pi$-orbit). The resulting metric space $(\widetilde{\mathcal{W}}, \delta_\square)$ is \emph{compact}. Consequently, every sequence of symmetric matrices $\{A_n\}_{n=1}^\infty$ with entries in a fixed interval $[m, M]$, after the affine normalization $a_{ij} \mapsto \frac{a_{ij} - m}{M - m} \in [0, 1]$, has a subsequence whose step realizations $W_{A_{n_k}}$ converge in $\delta_\square$ to some $W \in \widetilde{\mathcal{W}}$. The full sequence need not converge.
\end{theorem}
\begin{proof}
Compactness is proved by Lov\'asz and Szegedy \cite{lovasz2007szemeredi}, via the weak (Frieze--Kannan) regularity lemma and a martingale argument; the existence of limit objects for convergent graph sequences is due to the same authors \cite{lovasz2006limits}. See also the monograph \cite{lovasz2012large}.
\end{proof}

\begin{example}[The Identity Matrix Continuum Collapse]
Consider the identity matrix $I_n \in \R^{n \times n}$. In linear algebra, $\rank(I_n) = n \to \infty$. However, its functional realization $W_{I_n}(x, y)$ is supported on $n$ diagonal squares of area $1/n^2$, giving total support measure $n \times (1/n^2) = 1/n$. For any measurable sets $S, T \subseteq [0, 1]$:
\begin{equation}
    \cutnorm{W_{I_n}} \le \iint_{[0,1]^2} W_{I_n}(x, y) \dif x \dif y = \frac{1}{n} \longrightarrow 0 \quad \text{as } n \to \infty.
\end{equation}
Thus, in the continuous cut metric, $I_n \to 0$. Discrete full rank vanishes into the null graphon because the continuous diagonal has Lebesgue measure zero.
\end{example}

% =========================================================================
\section{High-Dimensional Tensors: Functional Geometry, Curvature, and Scaled Invariants}
\label{sec:tensors}
% =========================================================================

When transitioning from 2-tensors (matrices) to higher-order tensors $\mathcal{T} \in \bigotimes_{\alpha=1}^k \R^{d_\alpha}$ ($k \ge 3$), classical multilinear algebra confronts severe algebraic obstructions: tensor rank computation is NP-hard (Hillar and Lim~\cite{hillar2013most}), and low-rank tensor varieties fail to be topologically closed. In this section, we develop the functional machinery that overcomes these bottlenecks and extends our theory to high-order continuum limits.

\subsection{Variational Realization and Resolution of the Border-Rank Pathology}
\label{subsec:tensors_border_rank}

Let $\mathcal{T} \in \bigotimes_{\alpha=1}^k \R^{d_\alpha}$ be a $k$-th order tensor. The canonical CP approximation seeks:
\begin{equation}
    \label{eq:cp_approx_problem}
    \inf_{\mathcal{X} \in \mathcal{S}_r} \|\mathcal{T} - \mathcal{X}\|_F, \quad \text{where } \mathcal{S}_r \coloneqq \left\{ \sum_{i=1}^r \lambda_i u_i^{(1)} \otimes \dots \otimes u_i^{(k)} \right\}.
\end{equation}

\begin{theorem}[Functional Well-Posedness over Compact Product Manifolds]
\label{thm:tensor_weierstrass}
Define the product of unit spheres $\mathcal{M} \coloneqq \prod_{\alpha=1}^k \Sph^{d_\alpha - 1}$, which is compact under the product Riemannian metric. The functional realization:
\begin{equation}
    \Phi_{\mathrm{multi}}(\mathcal{T})(u^{(1)}, \dots, u^{(k)}) \coloneqq \sum_{j_1, \dots, j_k} \mathcal{T}_{j_1 \dots j_k} u_{j_1}^{(1)} \dots u_{j_k}^{(k)}
\end{equation}
is smooth: $\Phi_{\mathrm{multi}}(\mathcal{T}) \in C^\infty(\mathcal{M})$. Consequently:
\begin{enumerate}[label=(\roman*)]
    \item The maximal multilinear singular value:
    \begin{equation}
        \sigma_{\max}(\mathcal{T}) \coloneqq \sup_{(u^{(1)}, \dots, u^{(k)}) \in \mathcal{M}} \Phi_{\mathrm{multi}}(\mathcal{T})(u^{(1)}, \dots, u^{(k)})
    \end{equation}
    is attained at some tuple $(u_*^{(1)}, \dots, u_*^{(k)}) \in \mathcal{M}$.
    \item The critical points on $\mathcal{M}$ satisfy the Euler--Lagrange equations:
    \begin{equation}
        \label{eq:tensor_critical}
        \mathcal{T}(u^{(1)}, \dots, u^{(\alpha-1)}, \cdot, u^{(\alpha+1)}, \dots, u^{(k)}) = \sigma u^{(\alpha)}, \quad \forall \alpha \in \{1, \dots, k\},
    \end{equation}
    which are Lim's singular value equations for tensors~\cite{lim2005singular}. For a symmetric tensor and $u^{(1)} = \dots = u^{(k)}$ they reduce to the Z-eigenvalue equations of Lim and Qi~\cite{lim2005singular,qi2005eigenvalues}.
\end{enumerate}
In particular, a best rank-one approximation $\sigma_{\max}\, u_*^{(1)} \otimes \dots \otimes u_*^{(k)}$ of $\mathcal{T}$ always exists.
\end{theorem}
\begin{proof}
Since each sphere $\Sph^{d_\alpha - 1}$ is compact, by Tychonoff's Theorem $\mathcal{M}$ is compact. The polynomial mapping $\Phi_{\mathrm{multi}}(\mathcal{T})$ is continuous on $\mathcal{M}$. By the Weierstrass Extreme Value Theorem, $\Phi_{\mathrm{multi}}(\mathcal{T})$ attains its supremum and infimum on $\mathcal{M}$. 

To establish (ii), form the Lagrangian on $\prod \R^{d_\alpha} \times \R^k$:
\begin{equation}
    \mathcal{L}(u^{(1)}, \dots, u^{(k)}, \mu_1, \dots, \mu_k) = \Phi_{\mathrm{multi}}(\mathcal{T})(u^{(1)}, \dots, u^{(k)}) - \frac{1}{2} \sum_{\alpha=1}^k \mu_\alpha (\|u^{(\alpha)}\|_2^2 - 1).
\end{equation}
Vanishing of the gradient with respect to $u^{(\alpha)}$ gives:
\begin{equation}
    \nabla_{u^{(\alpha)}} \Phi = \mu_\alpha u^{(\alpha)}.
\end{equation}
Taking inner products with $u^{(\alpha)}$ and using multilinearity:
\begin{equation}
    \inner{u^{(\alpha)}}{\nabla_{u^{(\alpha)}} \Phi} = \Phi_{\mathrm{multi}}(\mathcal{T}) = \mu_\alpha \|u^{(\alpha)}\|_2^2 = \mu_\alpha.
\end{equation}
Since $\|u^{(\alpha)}\|_2 = 1$ for all $\alpha$, we must have $\mu_1 = \dots = \mu_k = \sigma$, so all singular equations share one parameter $\sigma$. For the rank-one statement, note that $\|\mathcal{T} - \lambda\, u^{(1)}\otimes\dots\otimes u^{(k)}\|_F^2 = \|\mathcal{T}\|_F^2 - 2\lambda\,\Phi_{\mathrm{multi}}(\mathcal{T})(u) + \lambda^2$ is minimized at $\lambda = \Phi_{\mathrm{multi}}(\mathcal{T})(u)$, which reduces the problem to maximizing $\Phi_{\mathrm{multi}}(\mathcal{T})^2$ over the compact set $\mathcal{M}$.
\end{proof}

\begin{remark}[Scope]
Compactness of $\mathcal{M}$ settles only the rank-one problem, whose feasible set is closed. For rank $r \ge 2$ the weights $\lambda_i$ in \eqref{eq:cp_approx_problem} range over $\R$. Restricting the factors to spheres does not bound them, and the non-closedness of $\mathcal{S}_r$ established by De Silva and Lim~\cite{desilva2008illposed} persists.
\end{remark}

\subsection{Continuum Limits of Dense Hyper-Networks: Hypergraphons}
\label{subsec:hypergraphons}

Just as symmetric matrices represent graphs, order-$k$ tensors $\mathcal{T} \in \bigotimes^k \R^n$ represent $k$-uniform hypergraphs (e.g., $k$-way multi-body interactions).

\begin{definition}[$k$-Hypergraphon and Multilinear Cut Norm]
A \emph{$k$-hypergraphon} is a symmetric measurable function $W \colon [0, 1]^k \to [0, 1]$. The \emph{Multilinear Cut Norm} of order $k$ is defined by:
\begin{equation}
    \label{eq:hypergraphon_cutnorm}
    \|W\|_{\square, k} \coloneqq \sup_{S_1, \dots, S_k \subseteq [0, 1]} \left| \int_{S_1 \times \dots \times S_k} W(x_1, \dots, x_k) \dif x_1 \dots \dif x_k \right|,
\end{equation}
where the supremum is taken over all measurable subsets $S_1, \dots, S_k \subseteq [0, 1]$.
\end{definition}

\begin{theorem}[Multilinear Operator Duality and Hypergraphon Compactness]
\label{thm:hypergraphon_compactness}
Let $\mathcal{W}_k$ be the space of $k$-hypergraphons endowed with the cut distance:
\begin{equation}
    \delta_{\square, k}(U, W) \coloneqq \inf_{\psi \in \bar{\Pi}} \|U - W^\psi\|_{\square, k},
\end{equation}
where $\bar{\Pi}$ is the group of measure-preserving bijections of $[0, 1]$. Then:
\begin{enumerate}[label=(\roman*)]
    \item The cut norm is equivalent to the $k$-multilinear operator norm from $L^\infty \to L^1$:
    \begin{equation}
        \|W\|_{\square, k} \le \|T_W\|_{L^\infty \times \dots \times L^\infty \to \R} \le 2^k \|W\|_{\square, k}.
    \end{equation}
    \item The quotient space $(\widetilde{\mathcal{W}}_k, \delta_{\square, k})$, in which hypergraphons at distance $0$ are identified, is \emph{compact}. Consequently, every sequence of symmetric tensors $\mathcal{T}_n \in \operatorname{Sym}^k(\R^n)$ with entries in a fixed interval $[m, M]$ has, after the affine normalization $t \mapsto \frac{t-m}{M-m}$ of the entries, a subsequence whose step realizations converge in $\delta_{\square,k}$. A non-symmetric tensor can first be replaced by its symmetrization, which changes the object under study.
\end{enumerate}
\end{theorem}
\begin{proof}
The proof generalizes Theorem~\ref{thm:grothendieck_cutnorm}. Decomposing each test function $u^{(\alpha)} \in L^\infty([0, 1])$ with $\|u^{(\alpha)}\|_\infty \le 1$ into positive and negative parts $u_+^{(\alpha)} - u_-^{(\alpha)}$, the multilinear integral splits into $2^k$ signed orthant integrals, each bounded by $\|W\|_{\square, k}$. The compactness of $\widetilde{\mathcal{W}}_k$ in this (product-set) cut distance is proved by Zhao~\cite{zhao2015hypergraph} via a weak regularity lemma; see also Gowers~\cite{gowers2007hypergraph} and Lov{\'a}sz~\cite{lovasz2012large}. For this cut norm the counting lemma fails for general $k$-uniform hypergraph densities, so the limit captures fewer invariants than in the graph case.
\end{proof}

\subsection{Multivariate Sobolev Regularization on Transformer Attention Tensors}
\label{subsec:transformer_tensors}

In modern deep learning architectures (Transformers), self-attention computes an activation tensor $\mathcal{A} \in \R^{B \times H \times N \times N}$, where $B$ is batch size, $H$ is number of attention heads, and $N$ is sequence length:
\begin{equation}
    \mathcal{A}_{b, h, i, j} = \operatorname{Softmax}\left( \frac{Q_{b, h, i, :} K_{b, h, j, :}^T}{\sqrt{d_k}} \right).
\end{equation}

\begin{proposition}[Sobolev Bound for Attention Realizations]
\label{thm:attention_sobolev}
Let $f_{\mathcal{A}_{b,h}}$ be the harmonic realization \eqref{eq:fourier_realization} of the attention slice $\mathcal{A}_{b, h, :, :}$, and fix $\lambda > 0$. Define the penalty, with the normalized measure $\dif\mu$ of \eqref{eq:fourier_realization},
\begin{equation}
    \mathcal{R}_{\mathrm{Attn}}(\mathcal{A}) \coloneqq \sum_{b=1}^B \sum_{h=1}^H \int_{\Tor^2} \left( |f_{\mathcal{A}_{b,h}}|^2 + \|\nabla f_{\mathcal{A}_{b,h}}\|_2^2 + \lambda |\Delta f_{\mathcal{A}_{b,h}}|^2 \right) \dif \mu,
\end{equation}
which is a sum of squared $H^2(\Tor^2)$ norms (with a $\lambda$-dependent weight on the second derivatives). Then:
\begin{enumerate}[label=(\roman*)]
    \item for every $b, h$,
    \begin{equation}
        \label{eq:attention_sup_bound}
        \sup_{(x, y) \in \Tor^2} |f_{\mathcal{A}_{b,h}}(x, y)| \le C_\lambda \sqrt{\mathcal{R}_{\mathrm{Attn}}(\mathcal{A})}, \qquad C_\lambda^2 \coloneqq \sum_{\mathbf{k} \in \Z^2} \frac{1}{1 + |\mathbf{k}|^2 + \lambda |\mathbf{k}|^4} < \infty;
    \end{equation}
    the constant $C_\lambda$ depends on $\lambda$ and tends to $\infty$ as $\lambda \to 0$;
    \item for every softmax attention matrix (entries $\ge 0$, rows summing to $1$), $f_{\mathcal{A}_{b,h}}(0,0) = \sum_{i,j} \mathcal{A}_{b,h,i,j} = N$. Hence $\sup |f_{\mathcal{A}_{b,h}}| \ge N$ and $\mathcal{R}_{\mathrm{Attn}}(\mathcal{A}) \ge N^2 / C_\lambda^2$.
\end{enumerate}
\end{proposition}
\begin{proof}
(i) Write $f = \sum_{\mathbf{k}} a_{\mathbf{k}} e^{i \mathbf{k}\cdot x}$ (finitely many terms). By Parseval for $\dif\mu$, the integrand of one slice integrates to $\sum_{\mathbf{k}} (1 + |\mathbf{k}|^2 + \lambda |\mathbf{k}|^4) |a_{\mathbf{k}}|^2$. By Cauchy--Schwarz,
\[
|f(x)| \le \sum_{\mathbf{k}} |a_{\mathbf{k}}| \le \Big( \sum_{\mathbf{k}} (1 + |\mathbf{k}|^2 + \lambda |\mathbf{k}|^4) |a_{\mathbf{k}}|^2 \Big)^{1/2} C_\lambda,
\]
and one slice is bounded by the sum over all slices. The series defining $C_\lambda^2$ converges because its terms are $\mathcal{O}(\lambda^{-1}|\mathbf{k}|^{-4})$. For $\lambda = 0$ the series $\sum_{\mathbf{k}} (1 + |\mathbf{k}|^2)^{-1}$ diverges, and by monotone convergence $C_\lambda \to \infty$ as $\lambda \to 0$. This is the Sobolev embedding $H^2(\Tor^2) \hookrightarrow L^\infty(\Tor^2)$, with nothing specific to attention. (ii) All exponentials equal $1$ at the origin.
\end{proof}

By (ii), the bound \eqref{eq:attention_sup_bound} cannot be uniform in the sequence length $N$: a penalty that stays bounded as $N$ grows is incompatible with softmax normalization under this realization. The statement concerns the attention field alone. It says nothing about downstream behaviour, such as hallucination, which depends on the whole decoding pipeline.

\subsection{Complexity of the High-Dimensional Tensor Morse Landscape}
\label{subsec:kac_rice_tensors}

\begin{theorem}[Complexity of the Isotropic Tensor Landscape]
\label{thm:kac_rice_tensors}
Let $d \ge 2$ and let $\mathcal{T} \in (\R^n)^{\otimes d}$ have independent standard Gaussian entries (no symmetry imposed), and let $f_{\mathcal{T}}(x) = \mathcal{T}(x, \dots, x)$ on $\Sph^{n-1}$. Equivalently, $f_{\mathcal{T}}$ is the form of the symmetric tensor $\operatorname{sym}\mathcal{T}$, whose entry indexed by a multiset $\mathbf{m}$ has variance $1/\operatorname{mult}(\mathbf{m})$, where $\operatorname{mult}(\mathbf{m})$ is the number of distinct orderings of $\mathbf{m}$. Then:
\begin{enumerate}[label=(\roman*)]
    \item $f_{\mathcal{T}}$ is a centered Gaussian field with covariance $\mathbb{E}[f_{\mathcal{T}}(x) f_{\mathcal{T}}(y)] = (x \cdot y)^d$, hence isotropic; this is the spherical pure $d$-spin model;
    \item for $d = 2$, $f_{\mathcal{T}}$ has exactly $2n$ critical points almost surely;
    \item for $d \ge 3$, the expected number of critical points grows exponentially, with the rate computed by Auffinger, Ben~Arous and \v{C}ern\'y~\cite{auffinger2013random}:
    \begin{equation}
        \label{eq:abc_complexity}
        \lim_{n \to \infty} \frac{1}{n} \log \mathbb{E}\big[|\operatorname{Crit}(f_{\mathcal{T}})|\big] = \frac{1}{2} \log(d-1).
    \end{equation}
    \item If instead the entries of a symmetric tensor $\mathcal{T} \in \operatorname{Sym}^d(\R^n)$ indexed by distinct multisets are independent with a common variance, then $f_{\mathcal{T}}$ is not isotropic for $n \ge 2$: $\operatorname{Var} f_{\mathcal{T}}(e_1) = 1$, whereas $\operatorname{Var} f_{\mathcal{T}}\big((e_1+e_2)/\sqrt{2}\big) = 2^{-d}\sum_{j=0}^d \binom{d}{j}^2 = 2^{-d}\binom{2d}{d}$, which is $\tfrac32$ for $d=2$ and $\tfrac52$ for $d=3$. Part (iii) does not apply to this model. What remains true for it, and for any law on $\operatorname{Sym}^d(\R^n)$ with a density, is the almost sure deterministic bound
    \begin{equation}
        \label{eq:cs_bound}
        |\operatorname{Crit}(f_{\mathcal{T}})| \le 2\, \frac{(d-1)^n - 1}{d-2} \quad (d \ge 3),
    \end{equation}
    so that $\limsup_n \frac1n \log \mathbb{E}|\operatorname{Crit}(f_{\mathcal{T}})| \le \log(d-1)$.
\end{enumerate}
\end{theorem}
\begin{proof}
(i) $f_{\mathcal{T}}(x) = \sum_{i_1,\dots,i_d} \mathcal{T}_{i_1\dots i_d} x_{i_1}\cdots x_{i_d}$ is a linear combination of independent standard Gaussians, so $\mathbb{E}[f_{\mathcal{T}}(x) f_{\mathcal{T}}(y)] = \sum_{i_1,\dots,i_d} x_{i_1} y_{i_1} \cdots x_{i_d} y_{i_d} = (x\cdot y)^d$. Grouping the $d!$-fold sums by multisets gives the stated variances of $\operatorname{sym}\mathcal{T}$.
(ii) For $d = 2$, $f_{\mathcal{T}}(x) = x^T S x$ with $S = \frac12(\mathcal{T} + \mathcal{T}^T)$, a GOE matrix, whose eigenvalues are almost surely distinct; Theorem~\ref{thm:morse_matrix} gives exactly $2n$ critical points.
(iii) The number of critical points does not change when $f_{\mathcal{T}}$ is multiplied by a positive constant or when the sphere is dilated. After these changes $f_{\mathcal{T}}$ becomes the Hamiltonian $H_{N,p}$ of~\cite{auffinger2013random}, with $N = n$ and $p = d$, on the sphere of radius $\sqrt{N}$, where the couplings are i.i.d.\ standard Gaussians indexed by ordered tuples. The total complexity of that model is $\frac12\log(p-1)$~\cite{auffinger2013random}. For orientation, the Kac--Rice formula (Adler and Taylor~\cite{adler2007random}) reads
\begin{equation}
    \mathbb{E}[|\operatorname{Crit}(f_{\mathcal{T}})|] = \int_{\Sph^{n-1}} \mathbb{E}\left[ |\det(\nabla^2 f_{\mathcal{T}}(x))| \;\middle|\; \nabla f_{\mathcal{T}}(x) = 0 \right] p_{\nabla f_{\mathcal{T}}(x)}(0) \dif \sigma(x).
\end{equation}
Differentiating the covariance $(x\cdot y)^d$ at $x = y$ shows the following. The tangential gradient has covariance $d\, I_{n-1}$. The tangential part $H_0$ of the ambient Hessian is uncorrelated with $f_{\mathcal{T}}(x)$ and with the gradient, and it is a GOE-type matrix with $\mathbb{E}[(H_0)_{ij}^2] = d(d-1)(1+\delta_{ij})$. By Euler's identity $x \cdot \nabla f_{\mathcal{T}}(x) = d\, f_{\mathcal{T}}(x)$, the Riemannian Hessian is
\begin{equation}
    \nabla^2 f_{\mathcal{T}}(x) = H_0 - d\, f_{\mathcal{T}}(x)\, I_{n-1}.
\end{equation}
For $d = 2$ this is the formula $2u^TAu - 2\lambda_i\|u\|^2$ of Theorem~\ref{thm:morse_matrix}. The exponential rate then comes from the large deviations of the GOE spectrum and an optimization over the energy level; this is carried out in~\cite{auffinger2013random}.
(iv) The variance of $f_{\mathcal{T}}(x) = \sum_{\mathbf{m}} \operatorname{mult}(\mathbf{m})\, \mathcal{T}_{\mathbf{m}} x^{\mathbf{m}}$ is $\sum_{\mathbf{m}} \operatorname{mult}(\mathbf{m})^2 x^{2\mathbf{m}}$. At $x = (e_1+e_2)/\sqrt2$ the multisets with $j$ copies of $1$ and $d-j$ copies of $2$ contribute $\binom{d}{j}^2 2^{-d}$. Critical points of $f_{\mathcal{T}}$ on the sphere are the real unit vectors $x$ with $\mathcal{T}(\cdot, x, \dots, x) = \mu x$, i.e.\ real eigenvectors. For a generic symmetric tensor the eigenvectors, counted in $\mathbb{P}^{n-1}(\C)$, are finite in number, and there are exactly $((d-1)^n - 1)/(d-2)$ of them (Cartwright and Sturmfels~\cite{cartwright2013number}). The non-generic tensors form a proper algebraic subset of $\operatorname{Sym}^d(\R^n)$, which has measure zero. Each real projective eigenvector gives two unit critical points $\pm x$, which proves \eqref{eq:cs_bound}. Letting $d \to 2$ in the bound recovers the count $2n$ of (ii).
\end{proof}

\subsection{Critical Scaling and Measure Concentration in Infinite Dimensions (\texorpdfstring{$d \to \infty$}{d -> infty})}
\label{subsec:measure_concentration}

A central open question in high-dimensional tensor analysis is identifying the exact normalization that prevents multilinear potentials from either diverging or vanishing as the mode dimension $d \to \infty$.

\begin{theorem}[Critical Multilinear Scaling and Sub-Gaussian Concentration]
\label{thm:critical_scaling}
Let $\mathcal{T} \in (\R^d)^{\otimes k}$ be a random Gaussian tensor whose entries are independent standard normal variables: $\mathcal{T}_{j_1 \dots j_k} \stackrel{\mathrm{i.i.d.}}{\sim} \mathcal{N}(0, 1)$. Define the normalized functional realization on unit spheres:
\begin{equation}
    \label{eq:renorm_tensor_form}
    \widehat{f}_{\mathcal{T}}(u^{(1)}, \dots, u^{(k)}) \coloneqq \frac{1}{\sqrt{d}} \sum_{j_1, \dots, j_k = 1}^d \mathcal{T}_{j_1 \dots j_k} u_{j_1}^{(1)} \dots u_{j_k}^{(k)}, \quad u^{(\alpha)} \in \Sph^{d-1}.
\end{equation}
Then:
\begin{enumerate}[label=(\roman*)]
    \item The normalization exponent $\alpha = 1/2$ is the unique critical scaling ensuring that the operator norm over unit spheres remains bounded and non-trivial:
    \begin{equation}
        \mathbb{E}\left[ \sup_{u^{(\alpha)} \in \Sph^{d-1}} |\widehat{f}_{\mathcal{T}}(u^{(1)}, \dots, u^{(k)})| \right] = \Theta(1) \quad \text{as } d \to \infty.
    \end{equation}
    \item Let $(u^{(1)}, \dots, u^{(k)})$ be random vectors on $\Sph^{d-1}$, with any joint law, independent of $\mathcal{T}$ (for instance i.i.d.\ uniform). Then, under the joint law of $(\mathcal{T}, u^{(1)}, \dots, u^{(k)})$, the value $\sqrt{d}\,\widehat{f}_{\mathcal{T}}(u^{(1)}, \dots, u^{(k)})$ is exactly standard normal, and therefore
    \begin{equation}
        \label{eq:concentration_inequality}
        \mathbb{P}\left( |\widehat{f}_{\mathcal{T}}(u^{(1)}, \dots, u^{(k)})| > \epsilon \right) \le 2 \exp\left( - \frac{d \, \epsilon^2}{2} \right), \qquad \epsilon > 0.
    \end{equation}
    The constant $\frac12$ in the exponent does not depend on $k$ and cannot be increased, since $\mathbb{P}(|Z| > s) = \exp(-\frac{s^2}{2}(1+o(1)))$ as $s \to \infty$ for $Z \sim \mathcal{N}(0,1)$.
\end{enumerate}
\end{theorem}

\begin{remark}[Spin-Glass Extensive Scaling vs Geometric Operator Norm]
In the statistical physics of spherical $p$-spin glasses ($p = k$), spin vectors are traditionally normalized to length $\|\sigma^{(\alpha)}\|_2 = \sqrt{d}$ so that each spin coordinate has variance $1$. In that convention, the multilinear sum pulls out a factor of $(\sqrt{d})^k = d^{k/2}$. The Hamiltonian is normalized by $d^{-(k-1)/2}$, yielding an extensive Hamiltonian $H_k(\sigma) = d^{-(k-1)/2} \sum \mathcal{T}_{j_1 \dots j_k} \sigma_{j_1}^{(1)} \dots \sigma_{j_k}^{(k)} = \sqrt{d} \sum \mathcal{T} u^{(1)} \dots u^{(k)}$. The intensive energy per particle $\frac{1}{d} H_k(\sigma) = \frac{1}{\sqrt{d}} \sum \mathcal{T} u^{(1)} \dots u^{(k)} = \widehat{f}_{\mathcal{T}}(u)$ is of order $\Theta(1)$, reconciling the physical Hamiltonian exponent $(k-1)/2$ with the geometric unit-sphere normalization factor $1/\sqrt{d}$.
\end{remark}

\begin{proof}
To prove (i), let $\mathcal{M} = \prod_{\alpha=1}^k \Sph^{d-1}$ be the product manifold of unit spheres. For an arbitrary fixed tuple $v \in \mathcal{M}$, the multilinear evaluation:
\begin{equation}
    X(v) \coloneqq \sum_{j_1, \dots, j_k = 1}^d \mathcal{T}_{j_1 \dots j_k} v_{j_1}^{(1)} \dots v_{j_k}^{(k)}
\end{equation}
is an exact linear combination of independent standard Gaussian random variables with variance:
\begin{equation}
    \operatorname{Var}(X(v)) = \sum_{j_1, \dots, j_k = 1}^d (v_{j_1}^{(1)})^2 \dots (v_{j_k}^{(k)})^2 = \prod_{\alpha=1}^k \|v^{(\alpha)}\|_2^2 = 1.
\end{equation}
Hence $X(v) \sim \mathcal{N}(0, 1)$ is identically a standard normal scalar for every $v \in \mathcal{M}$. Let $\mathcal{N}_\epsilon$ be an optimal $\epsilon$-net of $\Sph^{d-1}$ with cardinality $|\mathcal{N}_\epsilon| \le (1 + 2/\epsilon)^d$. The product net $\mathcal{N} \coloneqq \prod_{\alpha=1}^k \mathcal{N}_\epsilon \subset \mathcal{M}$ satisfies $|\mathcal{N}| \le (1 + 2/\epsilon)^{k d}$. By the standard Gaussian maximal inequality over the finite net $\mathcal{N}$:
\begin{equation}
    \mathbb{E}\left[ \sup_{v \in \mathcal{N}} |X(v)| \right] \le \sqrt{2 \log(2 |\mathcal{N}|)} \le \sqrt{2 k d \log(1 + 2/\epsilon) + 2\log 2} = C(k, \epsilon) \sqrt{d}.
\end{equation}
By multilinear approximation on $\mathcal{M}$, choosing $\epsilon = \frac{1}{2k}$ guarantees that $\sup_{u \in \mathcal{M}} |X(u)| \le \frac{1}{1 - k\epsilon} \sup_{v \in \mathcal{N}} |X(v)| \le 2 \sup_{v \in \mathcal{N}} |X(v)|$. Therefore:
\begin{equation}
    \mathbb{E}\left[ \sup_{u \in \mathcal{M}} |X(u)| \right] \le 2 C(k) \sqrt{d} = \mathcal{O}(\sqrt{d}).
\end{equation}
Conversely, restricting $u^{(1)} = \dots = u^{(k-1)} = e_1$ yields the Gaussian contraction $\sum_{j_k} \mathcal{T}_{1 \dots 1 j_k} u_{j_k}^{(k)}$. Its supremum over $u^{(k)} \in \Sph^{d-1}$ is the Euclidean norm of the standard Gaussian vector $Z \coloneqq (\mathcal{T}_{1 \dots 1 1}, \dots, \mathcal{T}_{1 \dots 1 d})^T \sim \mathcal{N}(0, I_d)$, which satisfies $\mathbb{E}[\|Z\|_2] = \sqrt{2} \frac{\Gamma((d+1)/2)}{\Gamma(d/2)} = \sqrt{d}(1 - \mathcal{O}(1/d)) = \Omega(\sqrt{d})$. Thus the unnormalized expected supremum is strictly $\Theta(\sqrt{d})$. Dividing by $\sqrt{d} = d^{1/2}$ yields $\mathbb{E}[\sup |\widehat{f}_{\mathcal{T}}|] = \Theta(1)$, establishing uniqueness of the scaling exponent $1/2$ on unit spheres.

To establish (ii), condition on $u = (u^{(1)}, \dots, u^{(k)})$. Since $u$ is independent of $\mathcal{T}$, the conditional law of $X(u) = \sqrt{d}\,\widehat{f}_{\mathcal{T}}(u)$ is that of $X(v)$ for the fixed tuple $v = u$, which is $\mathcal{N}(0,1)$ by the variance computation above. The conditional law does not depend on $u$, so $X(u) \sim \mathcal{N}(0,1)$ unconditionally. Hence $\mathbb{P}(|\widehat{f}_{\mathcal{T}}(u)| > \epsilon) = \mathbb{P}(|Z| > \epsilon\sqrt{d}) \le 2 e^{-d\epsilon^2/2}$ by the standard Gaussian tail bound $\mathbb{P}(Z > s) \le e^{-s^2/2}$ for $s \ge 0$.
\end{proof}

% =========================================================================
\section{Deep Implications Across the Mathematical Sciences}
\label{sec:implications}
% =========================================================================

We now discuss five application areas. In three of them (Sections~\ref{subsec:imp_networks}, \ref{subsec:imp_tensors}, and \ref{subsec:imp_imaging}) the relevant mathematics is established. It consists largely of known results (Lov\'asz--Szegedy graph limits, compactness for rank one, and the Fleming--Rishel coarea formula for total variation) viewed through the realization framework. In the other two (Sections~\ref{subsec:imp_dl} and \ref{subsec:imp_physics}) we offer proposals and conjectures whose consequences remain open, as noted in each subsection.

% -------------------------------------------------------------------------
\subsection{Deep Learning: Sobolev Weight Regularization and Zero-Shot Width Scaling}
\label{subsec:imp_dl}
In standard deep neural network training, weights $W \in \R^{d_{\mathrm{out}} \times d_{\mathrm{in}}}$ are conventionally regularized via coordinate-blind penalties: Weight Decay $\|W\|_F^2$ or Spectral Norm constraints $\sigma_{\max}(W) \le 1$. By Theorem~\ref{thm:dirichlet_blindness}, two weight matrices with identical spectral profiles can exhibit a factor of $\Theta(n^2)$ difference in their Dirichlet spatial energy $\mathcal{E}(f_W)$. 

We propose, as a hypothesis to be tested empirically, that this spectral blindness may leave networks susceptible to high-frequency perturbations. We propose the \emph{Sobolev Regularized Loss}:
\begin{equation}
    \label{eq:sobolev_loss}
    \mathcal{L}_{\mathrm{Sobolev}}(\theta) \coloneqq \mathcal{L}_{\mathrm{task}}(\theta) + \sum_{\ell=1}^L \lambda_\ell \|\nabla f_{W_\ell}\|_{L^2(\Tor^2)}^2.
\end{equation}
Penalizing the functional Dirichlet energy pushes the weights toward lower frequency indices. Two caveats apply. First, hidden units of a layer can be permuted without changing the network function, while this penalty is not permutation-invariant, so it depends on an arbitrary ordering of neurons unless that ordering carries meaning (e.g.\ spatial layouts in convolutional or positional structures). Second, we do not prove any bound on Lipschitz stability or expressive capacity; both are empirical questions.

Theorem~\ref{thm:graphon_compactness} also suggests a protocol for \emph{width scaling}, which we state as a proposal. For symmetric weight matrices $W_n \in \R^{n \times n}$ with entries in a fixed interval, the theorem gives only a \emph{subsequence} converging in $\delta_\square$ to some limit $\mathcal{W}_\infty \in \widetilde{\mathcal{W}}$. It does not say that the matrices produced by training at increasing widths converge, and it gives no rate. Non-symmetric weight matrices need the analogous theory for kernels in two variables that are not required to be symmetric. The limit is an equivalence class of functions defined almost everywhere, so point values $\mathcal{W}_\infty(\frac{i}{n}, \frac{j}{n})$ are not defined. A width-$n$ matrix can instead be obtained from a chosen representative by averaging over the cells $I_i \times I_j$, i.e.\ by the entries $n^2 \int_{I_i \times I_j} \mathcal{W}_\infty$. The proposal is to train at a moderate width (say $n = 10^3$), pass to a limit kernel, and read off matrices at a much larger width (say $n = 10^6$) in this way. Whether trained weights at moderate width are close in $\delta_\square$ to a common limit, and whether the transferred network performs well, are empirical questions not addressed by the theorem.

% -------------------------------------------------------------------------
\subsection{Theoretical Physics: Spin Glass Transitions and Parisi Symmetry Breaking}
\label{subsec:imp_physics}
In the statistical physics of disordered systems, the $p$-spin spherical spin glass is defined by the Hamiltonian $H_n(\sigma) = \sum_{i_1, \dots, i_p} J_{i_1 \dots i_p} \sigma_{i_1} \dots \sigma_{i_p}$ on the spin sphere $\sum \sigma_i^2 = n$. Determining the ground state energy and partition function $Z = \int e^{-\beta H_n} \dif \sigma$ via discrete tensor algebra is NP-hard.

Our functional realization on $\Sph^{n-1}$ (Theorem~\ref{thm:morse_matrix} and Theorem~\ref{thm:kac_rice_tensors}) places the spin glass landscape in the setting of Morse theory. For $p = 2$ there are $2n$ critical points, while for $p \ge 3$ the expected number grows like $\exp(\frac{n}{2}\log(p-1))$ \cite{auffinger2013random}. Auffinger, Ben~Arous, and \v{C}ern\'y also compute the complexity of critical points of each index below a given energy. They use it to identify the ground state energy and a threshold energy $-E_\infty$, below which the critical points of each fixed index are concentrated. This complexity result motivates a structural link between Giorgio Parisi's Replica Symmetry Breaking (RSB) and the proliferation of disconnected components in the sub-level sets
\begin{equation}
    M_E \coloneqq \left\{ \sigma \in \Sph^{n-1} \colon f_J(\sigma) \le E \right\}.
\end{equation}
We emphasize that identifying RSB with the onset of large zeroth Betti numbers $\beta_0(M_E) \gg 1$ is a structural conjecture rather than a consequence proved in this treatise: the critical-point count furnished by Theorem~\ref{thm:kac_rice_tensors} bounds the number of stationary points but does not by itself determine the connectivity of sub-level sets, which requires an independent Betti-number computation. We record this correspondence as an open problem in Section~\ref{sec:conclusion}.

% -------------------------------------------------------------------------
\subsection{Complex Networks: Size-Independent Metrics on Massive Graphs}
\label{subsec:imp_networks}
When comparing real-world complex networks of disparate dimensions (e.g., biological connectomes or internet topologies across developmental epochs, where $|V_1| = 10^5$ and $|V_2| = 10^8$), discrete matrix algebra provides no direct comparison: the difference $A_1 - A_2$ is undefined because the underlying vector spaces are non-isomorphic.

Under our step realization $\Phi_{\mathrm{step}}$, every adjacency matrix $A$ becomes a function $W_A \colon [0, 1]^2 \to [0, 1]$. By the graph-limit theory of Lov\'asz and Szegedy (Theorem~\ref{thm:graphon_compactness}), the cut distance $\delta_\square(W_{A_1}, W_{A_2})$ places networks of \emph{arbitrary, unequal sizes} in one compact metric space $(\widetilde{\mathcal{W}}, \delta_\square)$. This makes possible:
\begin{enumerate}[label=(\roman*)]
    \item Tracking a growing network through the curve $t \mapsto [W_{A(t)}] \in \widetilde{\mathcal{W}}$ and measuring its continuity in $\delta_\square$ (the quotient space has no linear structure, so time derivatives require working with representatives in $L^\infty([0,1]^2)$);
    \item Statistical hypothesis testing comparing graph families across different sampling scales;
    \item Continuous epidemic and consensus dynamics modeled via integro-differential equations rather than discrete master equations.
\end{enumerate}

% -------------------------------------------------------------------------
\subsection{Computational Multilinear Algebra: Circumventing the De Silva--Lim Pathology}
\label{subsec:imp_tensors}
In multilinear data processing (signal processing, hyper-spectral imaging, quantum chemistry), computing the best low-rank tensor approximation under the CP format:
\begin{equation}
    \min_{\rank(\mathcal{X}) \le r} \|\mathcal{T} - \mathcal{X}\|_F
\end{equation}
routinely encounters severe numerical instabilities. De Silva and Lim~\cite{desilva2008illposed} proved that for tensors of order $k \ge 3$, the variety $\mathcal{S}_r \coloneqq \{\mathcal{X} \colon \rank(\mathcal{X}) \le r\}$ is \emph{not topologically closed} in the Euclidean topology. Consequently, when the infimum is not attained, minimizing sequences must have factor weights that diverge to $\pm\infty$ while partially cancelling (degeneracy). Standard algorithms such as Alternating Least Squares (CP-ALS) can then show degenerate iterates or long plateaus (``swamps'').

The functional framework handles the rank-one case. Realizing the tensor as a multilinear form $\Phi(\mathcal{T})$ on the compact manifold $\mathcal{M} = \Sph^{n_1-1} \times \dots \times \Sph^{n_k-1}$ turns the best rank-one approximation into the maximization of a smooth function on a compact set, so a maximizer exists by the Weierstrass theorem. Lim's singular-vector equations~\cite{lim2005singular}, Eq.~\eqref{eq:tensor_critical}, characterize it. This is a classical fact: the set of rank-one tensors is closed. It does \emph{not} circumvent the De Silva--Lim pathology for $r \ge 2$ (see the remark after Theorem~\ref{thm:tensor_weierstrass}). Well-posed alternatives for higher rank require additional constraints, such as bounds on the weights or orthogonality or angular separation of the factors, or a change of format (e.g.\ Tucker or tensor-train formats, whose sets of tensors of bounded multilinear or TT rank are closed).

% -------------------------------------------------------------------------
\subsection{Biomedical Imaging: Hausdorff Perimeters and Minimal Surface Level Curves}
\label{subsec:imp_imaging}
In medical tomography (MRI, CT) and computer vision, digital images are discrete matrices of pixel intensities $A \in \R^{M \times N}$. The standard Frobenius norm $\|A\|_F^2 = \sum a_{ij}^2$ penalizes image energy diffusively, failing to distinguish between smooth background gradients and sharp organ boundaries.

By Theorem~\ref{thm:coarea_linkage}, the Total Variation $\TV(W_A)$ of the matrix step realization is mathematically identical to the \emph{integrated 1-dimensional Hausdorff measure} (arc-length) of the essential level curves:
\begin{equation}
    \TV(W_A) = \int_{-\infty}^\infty \mathcal{H}^1(\partial^* \{W_A > t\}) \dif t.
\end{equation}
This is the classical geometric reading of Total Variation regularization (as in the Rudin--Osher--Fatemi model~\cite{rudin1992nonlinear}): penalizing $\TV$ penalizes the integrated perimeter of the level sets. Matrices with identical pixel histograms can have very different level-set perimeters, so $\TV(W_A)$ is a natural quantitative descriptor of boundary complexity. Its diagnostic value for specific imaging tasks is an empirical question.

% =========================================================================
\section{Synthesis and Comparison Table}
\label{sec:synthesis}
% =========================================================================

Table~\ref{tab:comparison} synthesizes the correspondence between classical algebraic matrix invariants and their emerging functional counterparts.

\begin{table}[htbp]
\centering
\footnotesize
\renewcommand{\arraystretch}{1.3}
\begin{tabularx}{\textwidth}{@{}l X X X@{}}
\toprule
\textbf{Entity} & \textbf{Algebraic Invariant} & \textbf{Functional Realization} & \textbf{Emerging Invariant / Impact} \\
\midrule
$A \in \R^{n \times n}$ & Spectrum $\sigma(A)$, $\|A\|_F$ & $f_A \in H^s(\Tor^2)$ & Dirichlet Energy $\mathcal{E}(f_A)$, Sobolev DL regularization. \\
$A \in \R^{n \times n}$ & Finite differences $\Delta A$ & $W_A \in \BV((0,1)^2)$ & Total Variation $\TV$, Level-set perimeter $\mathcal{H}^1(\partial^* E_t)$. \\
$A \in \operatorname{Sym}_n$ & Spectrum $\lambda_1 \le \dots \le \lambda_n$ & $f_A \in C^\infty(\Sph^{n-1})$ & Morse indices $\gamma(v_i) = i-1$, Euler characteristic $\chi$. \\
$A$ ($n \to \infty$) & Rank, spectral radius $\rho$ & Graphon $W_A \in \widetilde{\mathcal{W}}$ & Cut norm $\cutnorm{W_A}$, Size-independent network metric. \\
$\mathcal{T} \in \operatorname{Sym}^d$ & Tensor rank, Z-eigenvalues & $f_{\mathcal{T}} \in C^\infty(\Sph^{n-1})$ & Critical-point complexity $\frac12\log(d-1)$ of the isotropic model ($d \ge 3$); well-posed rank-one approximation. \\
\bottomrule
\end{tabularx}
\caption{Systematic bridge between discrete multilinear algebra and emerging functional-geometric invariants.}
\label{tab:comparison}
\end{table}

% =========================================================================
\section{Conclusion and Open Problems}
\label{sec:conclusion}
% =========================================================================

We have formulated a comprehensive theoretical framework establishing that matrices and higher-order tensors are not merely discrete algebraic arrays, but discrete projections of rich continuous functional landscapes. Through Functional Realization Operators $\Phi$, we showed that properties traditionally considered external or non-algebraic---such as spatial smoothness, total variation, perimeter of level curves, Morse critical topologies, and infinite-dimensional graphon limits---arise naturally. The Dirichlet, total variation, Morse and cut-norm statements are proved here or are classical results with references. The attention bound (Proposition~\ref{thm:attention_sobolev}) is the Sobolev embedding and is not uniform in the sequence length. The complexity rate of Theorem~\ref{thm:kac_rice_tensors} is imported from Auffinger, Ben~Arous and \v{C}ern\'y.

Several fundamental open problems invite future research:
\begin{enumerate}
    \item \textbf{Optimal Realization Bounds:} For a given matrix $A$, what is the infimum Dirichlet energy $\inf_{\pi \in \mathcal{S}_n} \mathcal{E}(f_{P_\pi A P_\pi^T})$ over all coordinate permutations, and can this combinatorial-functional optimization problem be relaxed to a tractable Riemannian flow?
    \item \textbf{Hypergraphon Curvature:} Can one define a Ricci curvature tensor directly on higher-order tensor realizations (hypergraphons) that governs the convergence of message-passing algorithms on dense hypergraphs?
    \item \textbf{Topological Data Analysis (TDA) of Tensor Landscapes:} How do persistent homology barcodes of the sublevel sets $\{f_{\mathcal{T}} \le c\}$ correlate with the numerical rank and tensor train (TT) decomposition ranks of $\mathcal{T}$?
    \item \textbf{Betti Numbers and Replica Symmetry Breaking:} Does the zeroth Betti number $\beta_0(M_E)$ of the sub-level sets $M_E = \{\sigma \in \Sph^{n-1} \colon f_J(\sigma) \le E\}$ of the $p$-spin Hamiltonian undergo a sharp transition at the Parisi RSB temperature $T_c$, and can this transition be established rigorously from the critical-point complexity of Theorem~\ref{thm:kac_rice_tensors} together with an explicit connectivity estimate, rather than asserted by analogy?
\end{enumerate}

% Shared declarations block, \input by every chapter before its bibliography.
% Keep in sync with the "Declarations" section of master_book_unified_quantum_gravity.tex.
\section*{Declarations}

\noindent\textbf{Funding.} This research was financed in part by the Coordena\c{c}\~ao de Aperfei\c{c}oamento de Pessoal de N\'ivel Superior -- Brasil (CAPES) -- Finance Code 001.

\medskip\noindent\textbf{Competing interests.} The author declares no competing interests.

\medskip\noindent\textbf{Use of generative AI tools.} Large language model assistants (Anthropic Claude; Google Gemini, used through the Antigravity agent environment) were used extensively in preparing this work: drafting and editing text and \LaTeX{} source; proposing, expanding and revising derivations and proofs under the author's direction; writing the Python verification scripts and the \textsc{Lean 4} files; searching and checking the literature and references; and adversarial auditing of proofs, numerical claims and code. These audits found errors, some of them originally introduced with AI assistance. The errors and their corrections are recorded in the repository file \texttt{WORKPLAN.md}. AI tools are not authors. The author chose the research questions and the overall framework, reviewed the material, and is responsible for the content, including any remaining errors.

\medskip\noindent\textbf{Verification status.} Results labelled Theorem, Proposition or Lemma are proved in the text or cited as classical; statements labelled Conjecture are not proved. The numerical scripts compare formulas of the text with independent computations and are checks, not proofs. The \textsc{Lean 4} modules in \texttt{formal\_proofs\_book/Book} are a naming skeleton and do not verify the mathematics of this chapter.

\medskip\noindent\textbf{Code and data availability.} Sources, numerical scripts and the audit log are available at
\begin{center}\url{https://github.com/reinaldomsilvafilho-netizen/quantum-gravity-lean4}.\end{center}
The monograph is archived on Zenodo, \href{https://doi.org/10.5281/zenodo.22290043}{doi:10.5281/zenodo.22290043}, under the CC~BY~4.0 license.


\begin{thebibliography}{99}

\bibitem{adler2007random}
R.~J. Adler and J.~E. Taylor,
\emph{Random Fields and Geometry},
Springer Monographs in Mathematics, Springer, New York, 2007,
ISBN: 978-0-387-48112-8,
\href{https://doi.org/10.1007/978-0-387-48116-6}{DOI: 10.1007/978-0-387-48116-6}.

\bibitem{ambrosio2000functions}
L.~Ambrosio, N.~Fusco, and D.~Pallara,
\emph{Functions of Bounded Variation and Free Discontinuity Problems},
Oxford Mathematical Monographs, Oxford University Press, Oxford, 2000,
ISBN: 978-0-19-850245-6,
\href{https://doi.org/10.1093/oso/9780198502456.001.0001}{DOI: 10.1093/oso/9780198502456.001.0001}.

\bibitem{auffinger2013random}
A.~Auffinger, G.~Ben Arous, and J.~{\v{C}}ern{\`y},
\emph{Random matrices and complexity of spin glasses},
Communications on Pure and Applied Mathematics, \textbf{66} (2013), no.~2, 165--201,
\href{https://doi.org/10.1002/cpa.21422}{DOI: 10.1002/cpa.21422}.

\bibitem{cartwright2013number}
D.~Cartwright and B.~Sturmfels,
\emph{The number of eigenvalues of a tensor},
Linear Algebra and its Applications, \textbf{438} (2013), no.~2, 942--952,
\href{https://doi.org/10.1016/j.laa.2011.05.040}{DOI: 10.1016/j.laa.2011.05.040}.

\bibitem{desilva2008illposed}
V.~De Silva and L.-H. Lim,
\emph{Tensor rank and the ill-posedness of the best low-rank approximation problem},
SIAM Journal on Matrix Analysis and Applications, \textbf{30} (2008), no.~3, 1084--1127,
\href{https://doi.org/10.1137/06066518X}{DOI: 10.1137/06066518X}.

\bibitem{evans2015measure}
L.~C. Evans and R.~F. Gariepy,
\emph{Measure Theory and Fine Properties of Functions},
Revised Edition, Textbooks in Mathematics, CRC Press, Boca Raton, FL, 2015,
ISBN: 978-1-4822-4238-6,
\href{https://doi.org/10.1201/b18333}{DOI: 10.1201/b18333}.

\bibitem{fleming1960integral}
W.~H. Fleming and R.~Rishel,
\emph{An integral formula for total gradient variation},
Archiv der Mathematik, \textbf{11} (1960), 218--222,
\href{https://doi.org/10.1007/BF01236935}{DOI: 10.1007/BF01236935}.

\bibitem{federer1969geometric}
H.~Federer,
\emph{Geometric Measure Theory},
Die Grundlehren der mathematischen Wissenschaften, Band 153, Springer-Verlag, New York, 1969,
ISBN: 978-3-540-60656-7,
\href{https://doi.org/10.1007/978-3-642-62010-2}{DOI: 10.1007/978-3-642-62010-2}.

\bibitem{zhao2015hypergraph}
Y.~Zhao,
\emph{Hypergraph limits: a regularity approach},
Random Structures \& Algorithms \textbf{47} (2015), no.~2, 205--226,
\href{https://doi.org/10.1002/rsa.20537}{DOI: 10.1002/rsa.20537}.

\bibitem{gowers2007hypergraph}
W.~T. Gowers,
\emph{Hypergraph regularity and the multidimensional Szemer{\'e}di theorem},
Annals of Mathematics, \textbf{166} (2007), no.~3, 897--946,
\href{https://doi.org/10.4007/annals.2007.166.897}{DOI: 10.4007/annals.2007.166.897}.

\bibitem{hillar2013most}
C.~J. Hillar and L.-H. Lim,
\emph{Most tensor problems are NP-hard},
Journal of the ACM, \textbf{60} (2013), no.~6, Article 45,
\href{https://doi.org/10.1145/2512329}{DOI: 10.1145/2512329}.

\bibitem{lovasz2006limits}
L.~Lov{\'a}sz and B.~Szegedy,
\emph{Limits of dense graph sequences},
Journal of Combinatorial Theory, Series B, \textbf{96} (2006), no.~6, 933--957,
\href{https://doi.org/10.1016/j.jctb.2006.05.002}{DOI: 10.1016/j.jctb.2006.05.002}.

\bibitem{lovasz2007szemeredi}
L.~Lov{\'a}sz and B.~Szegedy,
\emph{Szemer{\'e}di's lemma for the analyst},
Geometric and Functional Analysis (GAFA), \textbf{17} (2007), no.~1, 252--270,
\href{https://doi.org/10.1007/s00039-007-0599-6}{DOI: 10.1007/s00039-007-0599-6}.

\bibitem{lovasz2012large}
L.~Lov{\'a}sz,
\emph{Large Networks and Graph Limits},
American Mathematical Society Colloquium Publications, Vol.~60, AMS, Providence, RI, 2012,
ISBN: 978-0-8218-9085-1,
\href{https://doi.org/10.1090/coll/060}{DOI: 10.1090/coll/060}.

\bibitem{lim2005singular}
L.-H. Lim,
\emph{Singular values and eigenvalues of tensors: a variational approach},
in \emph{1st IEEE International Workshop on Computational Advances in Multi-Sensor Adaptive Processing} (CAMSAP '05), 2005, pp.~129--132,
\href{https://doi.org/10.1109/CAMAP.2005.1574201}{DOI: 10.1109/CAMAP.2005.1574201}.

\bibitem{milnor1963morse}
J.~Milnor,
\emph{Morse Theory},
Annals of Mathematics Studies, No.~51, Princeton University Press, Princeton, NJ, 1963,
ISBN: 978-0-691-08008-6,
\href{https://doi.org/10.1515/9781400881802}{DOI: 10.1515/9781400881802}.

\bibitem{qi2005eigenvalues}
L.~Qi,
\emph{Eigenvalues of a real supersymmetric tensor},
Journal of Symbolic Computation, \textbf{40} (2005), no.~6, 1302--1324,
\href{https://doi.org/10.1016/j.jsc.2005.05.007}{DOI: 10.1016/j.jsc.2005.05.007}.

\bibitem{rudin1992nonlinear}
L.~I. Rudin, S.~Osher, and E.~Fatemi,
\emph{Nonlinear total variation based noise removal algorithms},
Physica D: Nonlinear Phenomena, \textbf{60} (1992), no.~1-4, 259--268,
\href{https://doi.org/10.1016/0167-2789(92)90242-F}{DOI: 10.1016/0167-2789(92)90242-F}.

\end{thebibliography}

\end{document}
